\clearpage
\begin{center}
\vspace*{3.5em}
{\LARGE\bfseries Álgebra de las magnitudes hipercomplejas\par}
\vspace{3.5em}
{\large Curso de la Prof. E. Noether\par}
{\large Semestre de invierno de 1929/30\par}
\vspace{1.5em}
{\large Redactado por el Prof. M. Deuring\par}
\vspace{4em}
{\Large\bfseries Índice\par}
\end{center}
\addcontentsline{toc}{section}{Álgebra de las magnitudes hipercomplejas}
\vspace{1em}
\begingroup
\small
\setlength{\parindent}{0pt}
\newcommand{\tocline}[2]{\noindent #1\dotfill #2\par}
\newcommand{\tocsec}[3]{\noindent\hspace*{1.2em}\S\,#1.\quad #2\dotfill #3\par}
\tocline{\textbf{Introducción}}{3}
\medskip
\tocline{\textbf{Capítulo I: Representaciones}}{3}
\tocsec{1}{Definición de la representación directa --- Definición de la representación recíproca}{3}
\tocsec{2}{Clases de representaciones}{4}
\tocsec{3}{Módulos de representación}{4}
\tocsec{4}{Relación entre módulos de representación y representaciones}{5}
\medskip
\tocline{\textbf{Capítulo II: Teoría de Galois de los cuerpos conmutativos}}{7}
\tocsec{5}{Extensión de coeficientes de los sistemas hipercomplejos}{7}
\tocsec{6}{Representaciones irreducibles de los sistemas conmutativos}{8}
\tocsec{7}{Isomorfismos de un cuerpo}{8}
\tocsec{8}{El cuerpo de descomposición}{9}
\tocsec{9}{Isomorfismos de $Z_1$ sobre los $Z_i$}{9}
\tocsec{10}{El grupo de Galois}{9}
\tocsec{11}{El teorema fundamental de la teoría de Galois}{9}
\tocsec{12}{Significado formal de las componentes $e_i$}{11}
\tocsec{13}{El teorema general de prolongación}{12}
\medskip
\tocline{\textbf{Capítulo III: Grupos abelianos}}{13}
\tocsec{14}{El anillo de grupo}{13}
\tocsec{15}{Los anillos de grupo de los grupos abelianos}{14}
\tocsec{16}{Las relaciones de caracteres}{15}
\tocsec{17}{La teoría de Galois de los grupos abelianos}{15}
\medskip
\tocline{\textbf{Capítulo IV: Anillos simples bilaterales}}{18}
\tocsec{18}{Un lema}{18}
\tocsec{19}{Representaciones de sistemas hipercomplejos simples bilaterales en cuerpos de extensión no conmutativos de sus cuerpos de coeficientes}{19}
\tocsec{20}{Cuerpos no conmutativos}{22}
\tocsec{21}{La teoría de Galois de los cuerpos no conmutativos}{26}
\clearpage
\tocline{\textbf{Capítulo V: Sistemas de factores}}{29}
\tocsec{22}{El grupo de los cuerpos con centro dado}{29}
\tocsec{23}{Sistemas de factores}{30}
\tocsec{24}{Multiplicación de sistemas de factores}{33}
\tocsec{25}{Representación normal de $\mathfrak R_r$ con subcuerpo conmutativo maximal de Galois}{39}
\tocsec{26}{Multiplicación de las representaciones cruzadas}{43}
\medskip
\tocline{\textbf{Capítulo VI: Teoría de los productos cruzados}}{44}
\tocsec{27}{Representaciones de los $\mathfrak R_r$ como productos cruzados}{44}
\tocsec{28}{Teorema del producto para sistemas de factores}{47}
\tocsec{29}{Teorema del género principal en el caso minimal}{50}
\tocsec{30}{Especialización a cuerpos de descomposición cíclicos}{50}
\tocsec{31}{Aplicaciones del caso especial cíclico}{52}
\endgroup
\clearpage

\section*{Introducción}

Este curso trata la teoría general de las representaciones de los anillos, surgida, por una parte, de la teoría de las representaciones de los grupos finitos y, por otra, de la teoría algebraica y aritmética de los sistemas de números hipercomplejos. Esta concepción más general de la teoría de representaciones presenta frente a la anterior (Frobenius, Schur) las ventajas siguientes:

Como no se considera el grupo, sino el anillo de grupo, o, más generalmente, un sistema hipercomplejo o un anillo arbitrario, puede emplearse la teoría de anillos e ideales, lo cual libera a la teoría de muchos cálculos poco transparentes. El uso de la teoría general de anillos e ideales no solo simplifica la teoría, sino que también la hace avanzar ---por ejemplo, en la cuestión de las relaciones entre un grupo y sus subgrupos--- y le abre nuevos campos de aplicación. Con ayuda de la teoría general de representaciones puede darse una fundamentación nueva y muy elegante de la teoría de Galois y ---lo que acaso sea aún más importante--- establecerse una teoría de Galois de los cuerpos no conmutativos.

Los conceptos fundamentales de la teoría de grupos, anillos e ideales, así como los teoremas generales de la teoría de representaciones, se encuentran en el trabajo de la señorita Noether, \emph{Magnitudes hipercomplejas y teoría de representaciones}, \emph{Mathematische Zeitschrift} 30, p.~641. Este trabajo, citado en lo que sigue como M.Z., debe considerarse complemento de la presente exposición. El mismo material de M.Z. se encuentra también en la redacción del curso de la señorita Noether sobre magnitudes hipercomplejas y caracteres de grupos del semestre de invierno de 1927/28.

Aquí se explican brevemente los conceptos de representación y módulo de representación para introducir el módulo de representación recíproca y las representaciones recíprocas, que solo aportan una novedad formal.

\section*{Capítulo I. Representaciones}

Sea $\mathfrak{o}$ un anillo ---el anillo que se representa--- y sea $\mathsf T$ otro anillo ---aquel en el cual se representa---.

\subsection*{§ 1. Definición de la representación directa}

Una \emph{representación directa} de grado $n$ de $\mathfrak{o}$ en $\mathsf T$ es un subanillo $\mathfrak D$ del anillo $\mathsf T_n$ de las matrices de grado $n$ con elementos de $\mathsf T$, sobre el cual $\mathfrak{o}$ se aplica homomórficamente como anillo. En símbolos,

\[
\mathfrak{o}\to\mathfrak D\subseteq\mathsf T_n.
\]

\subsection*{Definición de la representación recíproca}

Una \emph{representación recíproca} de grado $n$ de $\mathfrak{o}$ en $\mathsf T^*$ es un subanillo $\mathfrak D^*$ del anillo $\mathsf T_n^*$ de las matrices de grado $n$ con elementos de $\mathsf T^*$, sobre el cual $\mathfrak{o}$ se aplica antihomomórficamente como anillo. Aquí \emph{antihomomórficamente} significa lo siguiente: si a los elementos $c,d$ de $\mathfrak{o}$ corresponden los elementos $c^*,d^*$ de $\mathfrak D^*$, entonces al elemento $cd$ de $\mathfrak{o}$ debe corresponder el elemento $d^*\cdot c^*$ de $\mathfrak D^*$, y al elemento $c+d$, el elemento $c^*+d^*$. En símbolos,

\[
\mathfrak{o}\mapsto\mathfrak D^*\subseteq\mathsf T_n^*.
\]

Como ejemplo del concepto de antihomomorfía, la aplicación que asigna a cada matriz de $n$ filas sobre un cuerpo conmutativo su transpuesta es un antiisomorfismo de anillos.

En general, dado un anillo cualquiera $\mathsf T$, puede construirse un anillo $\mathsf T^*$ antiisomorfo a él del modo siguiente:

A cada elemento $c$ de $\mathsf T$ se le asigna un símbolo $c^*$ y se define la adición de estos símbolos mediante $c^*+d^*=(c+d)^*$, y la multiplicación mediante

\[
c^*\cdot d^*=(dc)^*.
\]

El conjunto de los símbolos $c^*$ es un anillo $\mathsf T^*$ antiisomorfo a $\mathsf T$.

\subsection*{§ 2. Clases de representaciones}

Una representación directa $\mathfrak D$ de $\mathfrak{o}$ en $\mathsf T$ puede transformarse en una representación isomorfa $\mathfrak D'$ mediante una transformación elemento a elemento con una matriz regular $P$; es decir, al elemento $c$ de $\mathfrak{o}$ se le asigna, en lugar de $C\in\mathfrak D$, el elemento $P^{-1}CP\in P^{-1}\mathfrak D P$.

Dos representaciones que pueden transformarse una en otra de este modo se llaman \emph{equivalentes}.

Todas las representaciones equivalentes a una representación dada forman una \emph{clase de representaciones} (c. r. d.).

La equivalencia de representaciones recíprocas se define de igual manera: dos representaciones recíprocas se llaman equivalentes cuando una se obtiene de la otra mediante una transformación con una matriz regular. Todas las representaciones equivalentes a una representación recíproca fija forman una \emph{clase de representaciones recíprocas}.

\subsection*{§ 3. Módulos de representación}

Definición del \emph{módulo de representación directa} (m. r. d.).

Un módulo de representación directa de $\mathfrak{o}$ respecto de $\mathsf T$ es todo bimódulo $\mathfrak{o}$--$\mathsf T$, $\mathfrak M$ (esto es, un grupo abeliano escrito aditivamente que es módulo izquierdo respecto de $\mathfrak{o}$, módulo derecho respecto de $\mathsf T$ y satisface la ley asociativa mixta: si $c\in\mathfrak{o}$, $m\in\mathfrak M$ y $\tau\in\mathsf T$, entonces $(cm)\tau=c(m\tau)$), que puede escribirse como suma directa de un número finito de módulos cíclicos sobre $\mathsf T$ ---con $\tau=0$ si $x_i\tau=0$---:

\[
\mathfrak M=x_1\cdot\mathsf T+x_2\cdot\mathsf T+\cdots+x_n\cdot\mathsf T.
\]

Definición del \emph{módulo de representación recíproca} (m. r. r.).

Un módulo de representación recíproca de $\mathfrak{o}$ respecto de $\mathsf T^*$ es todo módulo izquierdo sobre $\mathfrak{o}$ y sobre $\mathsf T^*$, $\mathfrak M^*$, que satisface la ley de conmutación (si $c\in\mathfrak{o}$, $m^*\in\mathfrak M^*$ y $\tau^*\in\mathsf T^*$, entonces $c(\tau^*m^*)=\tau^*(cm^*)$) y que puede escribirse como suma directa de un número finito de módulos simples sobre $\mathsf T^*$ ---con $\tau^*=0$ si $\tau^*x_i^*=0$---:

\[
\mathfrak M^*=\mathsf T^*\cdot x_1^*+\cdots+\mathsf T^*\cdot x_n^*.
\]

\subsection*{§ 4. Relación entre los módulos de representación y las representaciones}

\textbf{1.} \emph{Todo módulo de representación determina de manera única una clase de representaciones (un m. r. d. determina una c. r. d.; un m. r. r., una c. r. r.).}

Esto se obtiene del modo siguiente.

La multiplicación del m. r. d. $\mathfrak M$ por un elemento $c$ de $\mathfrak{o}$ es un endomorfismo de $\mathfrak M$ que, en virtud de la ley asociativa $(cm)\tau=c(m\tau)$, admite los elementos de $\mathsf T$ como operadores. Es, por tanto, una \emph{transformación lineal} de $\mathfrak M$ en sí mismo; es decir, la aplicación $m\mapsto cm=m'$ tiene la propiedad

\[
\text{si }y_i\mapsto y_i'(=cy_i),\text{ entonces también }
\sum y_i\tau_i\mapsto\sum y_i'\tau_i\left(=c\sum y_i\tau_i\right).
\]

Definamos la suma $\mathfrak T_1+\mathfrak T_2$ de dos transformaciones lineales

\[
\mathfrak T_1:m\mapsto m'(=c_1m),\qquad
\mathfrak T_2:m\mapsto m''(=c_2m)
\]

como la transformación $m\mapsto m'+m''=(c_1+c_2)m$, y su producto $\mathfrak T_1\mathfrak T_2$ como la transformación compuesta $m\mapsto(m'')'=c_1c_2m$. El conjunto de las distintas transformaciones lineales de $\mathfrak M$ engendradas por los elementos de $\mathfrak{o}$ es entonces una imagen homomorfa de anillos $\overline{\mathfrak D}$ de $\mathfrak{o}$.

Elijamos ahora una base cualquiera $x_1,\ldots,x_n$ de $\mathfrak M$. La transformación lineal engendrada por $c$ queda determinada cuando se conoce la transformación de los elementos de la base, esto es, cuando se conoce la matriz $(\gamma_{ij})=C$ del sistema

\[
x'_j=cx_j=\sum_i x_i\gamma_{ij},
\]

pues para cualquier otro elemento $x=\sum_jx_j\tau_j$ se tiene simplemente $x'=\sum_jx'_j\tau_j$. Estas matrices $C$ están en correspondencia biunívoca con las distintas transformaciones lineales y forman, como se ve inmediatamente, un anillo isomorfo a $\overline{\mathfrak D}$; constituyen, por tanto, una imagen homomorfa de anillos de $\mathfrak{o}$, es decir, una representación de grado $n$ de $\mathfrak{o}$.

El isomorfismo entre $\overline{\mathfrak D}$ y la representación se expresa en las fórmulas siguientes. Si a la transformación lineal engendrada por $c\in\mathfrak{o}$ corresponde la matriz $C$, de modo que

\[
c(x_1,\ldots,x_n)=(x_1,\ldots,x_n)C,
\]

y análogamente a $d\in\mathfrak{o}$ corresponde $D\in\mathsf T_n$, de modo que

\[
d(x_1,\ldots,x_n)=(x_1,\ldots,x_n)D,
\]

entonces

\[
\begin{aligned}
(c+d)(x_1,\ldots,x_n)
&=(cx_1+dx_1,\ldots,cx_n+dx_n)\\
&=(x_1,\ldots,x_n)(C+D),
\end{aligned}
\]

y

\[
\begin{aligned}
cd(x_1,\ldots,x_n)
&=c(dx_1,\ldots,dx_n)=c((x_1,\ldots,x_n)D)\\
&=(cx_1,\ldots,cx_n)D=(x_1,\ldots,x_n)CD.
\end{aligned}
\]

Así, al realizar mediante una base determinada $x_\nu$ de $\mathfrak M$ las transformaciones lineales de $\mathfrak M$ engendradas por $\mathfrak{o}$, se obtiene una representación de grado $n$ de $\mathfrak{o}$ en $\mathsf T$.

Si $(y_1,\ldots,y_n)$ es otra base de $\mathfrak M$, se obtiene de $x_1,\ldots,x_n$ por multiplicación por una matriz regular:

\[
(y_1,\ldots,y_n)=(x_1,\ldots,x_n)P.
\]

Si una transformación lineal de $\mathfrak M$ se realiza en la base $x_\nu$ mediante $C$, en la base $y_\nu$ se realiza mediante $P^{-1}CP$, pues de

\[
c(x_1,\ldots,x_n)=(x_1,\ldots,x_n)C
\]

se sigue

\[
\begin{aligned}
c(y_1,\ldots,y_n)
&=c(x_1,\ldots,x_n)P=((x_1,\ldots,x_n)C)P\\
&=(x_1,\ldots,x_n)CP=((y_1,\ldots,y_n)P^{-1})CP\\
&=(y_1,\ldots,y_n)P^{-1}CP.
\end{aligned}
\]

De aquí se sigue inmediatamente que $\mathfrak M$ produce todas las representaciones de una clase de representaciones cuando se utilizan todas sus bases para realizar las transformaciones lineales inducidas por $\mathfrak{o}$.

\textbf{2.} \emph{Toda clase de representaciones es engendrada por un módulo de representación de la manera indicada en 1.} Nos restringimos de nuevo a las representaciones directas.

Demostración. Tomemos una representación determinada de la clase, dada por la correspondencia homomorfa $c\mapsto C$. Si $n$ es su grado, formamos con $n$ indeterminadas $x_1,\ldots,x_n$ el módulo derecho sobre $\mathsf T$

\[
\mathfrak M=x_1\mathsf T+\cdots+x_n\mathsf T,
\]

con las reglas usuales de cálculo

\[
\sum x_i\tau_i+\sum x_i\bar\tau_i=\sum x_i(\tau_i+\bar\tau_i),
\qquad
\sum x_i\tau_i=\sum x_i\bar\tau_i\iff \tau_i=\bar\tau_i\text{ para todo }i,
\]

y $(\sum x_i\tau_i)\varrho=\sum x_i\tau_i\varrho$. Lo convertimos en módulo izquierdo sobre $\mathfrak{o}$ imponiendo

\[
c\sum x_i\tau_i=\sum x_i'\tau_i
\]

cuando $(x_1',\ldots,x_n')=(x_1,\ldots,x_n)C$, donde $C$ es la matriz de la representación asignada al elemento $c$ de $\mathfrak{o}$.

Con esta definición $\mathfrak M$ se convierte en un bimódulo $\mathfrak{o}$--$\mathsf T$. En efecto, de las relaciones de homomorfía

\[
c+d\mapsto C+D,\qquad cd\mapsto CD
\]

se siguen

\[
(c+d)x_i=cx_i+dx_i,
\]

\[
cd\,x_i=c(dx_i),\qquad
c\left(\sum x_i\tau_i+\sum x_i\bar\tau_i\right)
=c\sum x_i\tau_i+c\sum x_i\bar\tau_i;
\]

estas son las propiedades de módulo izquierdo. Además,

\[
c\left(\sum x_i\tau_i\varrho\right)=\left(c\sum x_i\tau_i\right)\varrho,
\]

que es la ley asociativa mixta.

Por tanto $\mathfrak M$ es un módulo de representación, y la representación elegida es engendrada por él del modo descrito en 1 al utilizar la base $x_1,\ldots,x_n$ para realizar sus transformaciones lineales. Las demás bases proporcionan las restantes representaciones de la clase.

Para las representaciones y módulos de representación recíprocos todo procede de la misma manera. El módulo

\[
\mathfrak M^*=\mathsf T^*x_1^*+\cdots+\mathsf T^*x_n^*
\]

engendra una representación recíproca de $\mathfrak{o}$ en $\mathsf T^*$ asignando al elemento $c$ de $\mathfrak{o}$ la matriz $C$ de la ecuación

\[
c\begin{pmatrix}x_1^*\\ \vdots\\ x_n^*\end{pmatrix}
=\begin{pmatrix}cx_1^*\\ \vdots\\ cx_n^*\end{pmatrix}
=C\begin{pmatrix}x_1^*\\ \vdots\\ x_n^*\end{pmatrix}.
\]

Las distintas bases vuelven a dar las distintas representaciones de una clase.

\section*{Capítulo II. Teoría de Galois de los cuerpos conmutativos}

La teoría de representaciones desarrollada hasta aquí basta para fundamentar la teoría de Galois y presenta, frente a la exposición usual, la ventaja de no tener que recurrir a sistemas especiales de elementos generadores ni a teoremas sobre polinomios en indeterminadas.\footnote{Para lo que sigue, compárese M.Z., §§15--20. M.Z., §21, contiene los primeros teoremas del apartado siguiente sobre teoría de Galois. Para la definición de los sistemas hipercomplejos, véase M.Z., §6.}

\subsection*{§ 5. Extensión de coeficientes de sistemas hipercomplejos}

Formulamos el siguiente teorema general, cuya demostración omitimos por ser casi inmediata. Sea

\[
\mathfrak{o}=c_1\mathsf P+\cdots+c_n\mathsf P
\]

un sistema hipercomplejo sobre el cuerpo conmutativo $\mathsf P$. Sea $\Omega$ un anillo de extensión de $\mathsf P$ cuyo centro contiene a $\mathsf P$. Definimos el \emph{sistema extendido} $\mathfrak{o}_\Omega$ como el conjunto de todas las sumas formales $\sum c_i\omega_i$, con $\omega_i\in\Omega$. Mediante las convenciones

\[
\begin{array}{l}
\sum c_i\omega_i=\sum c_i\omega_i'\quad\text{si y solo si }\omega_i=\omega_i'\text{ para todo }i;\\
\sum c_i\omega_i+\sum c_i\omega_i'=\sum c_i(\omega_i+\omega_i');\\
\omega\sum c_i\omega_i=\sum c_i\omega\omega_i;\qquad
\left(\sum c_i\omega_i\right)\omega=\sum c_i\omega_i\omega;\\
\sum c_i\omega_i=\sum\omega_i c_i,
\end{array}
\]

se convierte primero en un módulo sobre $\Omega$ conectado conmutativamente con $\Omega$; y mediante la convención adicional

\[
\left(\sum_i c_i\omega_i\right)\left(\sum_j c_j\omega_j\right)
=\sum_i\sum_j c_ic_j\omega_i\omega_j
\]

---donde $c_ic_j$ ya está definido como elemento de $\mathfrak{o}$--- se convierte en un anillo de extensión de $\mathfrak{o}$. El sistema $\mathfrak{o}_\Omega$ tiene rango $n$ sobre $\Omega$ y puede escribirse en la forma

\[
\mathfrak{o}_\Omega=c_1\Omega+\cdots+c_n\Omega.
\]

La definición de $\mathfrak{o}_\Omega$ es independiente de la base $c_1,\ldots,c_n$. Si $\Omega$ es un cuerpo conmutativo, $\mathfrak{o}_\Omega$ es un sistema hipercomplejo sobre $\Omega$.

\subsection*{§ 6. Las representaciones irreducibles de los sistemas conmutativos}

Sea $\mathfrak Z=z_1\mathsf P+\cdots+z_n\mathsf P$ un sistema conmutativo sobre $\mathsf P$. Queremos determinar todas las representaciones irreducibles de $\mathfrak Z$ en un cuerpo de extensión algebraicamente cerrado $\Omega$ de $\mathsf P$. Toda representación de $\mathfrak Z$ en $\Omega$ induce una representación de

\[
\mathfrak Z_\Omega=z_1\Omega+\cdots+z_n\Omega,
\]

pues las representaciones de los elementos de base ya contenidos en $\mathsf P$ bastan para determinar la representación entera; recíprocamente, toda representación de $\mathfrak Z_\Omega$ contiene una representación de $\mathfrak Z$. Podemos, por tanto, buscar las representaciones de $\mathfrak Z_\Omega$. Según M.Z., §20, estas están contenidas en la representación regular de $\mathfrak Z_\Omega$, esto es, en la representación que se obtiene empleando $\mathfrak Z_\Omega$ mismo como módulo de representación. Basta incluso restringirse al módulo de representación $\mathfrak Z_\Omega/\mathfrak C$, donde $\mathfrak C$ es el radical de $\mathfrak Z_\Omega$ (véase M.Z., §19).

Ahora bien, $\mathfrak Z_\Omega/\mathfrak C$ es un sistema completamente reducible y, por tanto, una suma directa de cuerpos de grado finito sobre $\Omega$. Como $\Omega$ es algebraicamente cerrado, esos sumandos son isomorfos a $\Omega$:

\[
\mathfrak Z_\Omega/\mathfrak C=e_1\Omega+\cdots+e_t\Omega.
\]

Los $e_\nu$ son las componentes de la unidad. Por consiguiente:

\emph{El sistema $\mathfrak Z_\Omega$ posee exactamente $t$ clases distintas de representaciones irreducibles en $\Omega$, todas de grado uno. Aquí $t$ es el rango de $\mathfrak Z_\Omega/\mathfrak C$ y, por tanto, $t\leq n$.}

Cada una de estas clases contiene una sola representación, pues transformar una representación de grado uno en un cuerpo conmutativo no la altera: $\pi^{-1}\alpha\pi=\alpha$.

Existen, por tanto, exactamente $t$ aplicaciones homomorfas distintas $\Theta_1,\ldots,\Theta_t$ de $\mathfrak Z_\Omega$ sobre subanillos de $\Omega$.

\subsection*{§ 7. Los isomorfismos de un cuerpo}

Supongamos ahora que $\mathfrak Z$ es un cuerpo. Todo homomorfismo $\Theta_i$ de $\mathfrak Z_\Omega$ sobre un subanillo de $\Omega$ debe aplicar isomórficamente $\mathfrak Z$ sobre un subcuerpo $Z_i$ de $\Omega$. En efecto, o bien $Z_i=0$ ---lo cual no ocurre aquí, pues a la unidad de $\mathfrak Z$ corresponde la unidad de $\Omega$---, o bien $Z_i$ es una imagen isomorfa de $\mathfrak Z$. Por consiguiente:

\emph{Si $\mathfrak Z$ es una extensión de grado finito $n$ del cuerpo $\mathsf P$, en $\Omega$ existen tantos cuerpos distintos $Z_i$ isomorfos a $\mathfrak Z$ como indique el rango $t$ de $\mathfrak Z_\Omega/\mathfrak C$. El número de los $Z_i$ es, pues, a lo sumo $n$.}

Se dice que $\mathfrak Z$ es de \emph{primera especie} sobre $\mathsf P$ cuando el número de los $Z_i$ es igual a $n$, y de \emph{segunda especie} cuando es menor que $n$.

El cuerpo $\mathfrak Z$ es de primera especie sobre $\mathsf P$ si y solo si $\mathfrak Z_\Omega$ es completamente reducible.

Aquí se aprecia la ventaja de esta fundamentación de la teoría de Galois frente a la habitual. En la teoría de Galois usual se consideran los isomorfismos de uno de los $Z_i$ sobre los demás; aquí se evita esa asimetría. Para demostrar que un cuerpo no posee más cuerpos conjugados que los indicados por su grado, suele emplearse un elemento primitivo. Aquí el hecho se sigue simplemente de que todas las representaciones irreducibles de $\mathfrak Z_\Omega$ están contenidas en la representación regular.

\subsection*{§ 8. El cuerpo de descomposición}

Los $Z_i$ son las representaciones irreducibles de $\mathfrak Z$ en $\Omega$. Por tanto, si un cuerpo $\Gamma$ intermedio entre $\mathsf P$ y $\Omega$,

\[
\mathsf P\subseteq\Gamma\subseteq\Omega,
\]

contiene todos los $Z_i$, entonces $\mathfrak Z_\Gamma/\mathfrak C$ debe descomponerse ya en $t$ cuerpos de grado uno sobre $\Gamma$. Recíprocamente, para que $\mathfrak Z_\Gamma/\mathfrak C$ se descomponga en $t$ cuerpos ---que serán entonces las representaciones irreducibles de grado uno de $\mathfrak Z_\Gamma$--- es necesario que los $Z_i$ estén contenidos en $\Gamma$. Si se llama \emph{cuerpo de descomposición} de $\mathfrak Z$ a un cuerpo $\Gamma$ tal que $\mathfrak Z_\Gamma/\mathfrak C$ se convierte en suma directa de $t$ cuerpos, obtenemos:

\emph{El cuerpo de Galois de los $Z_i$ isomorfos a $\mathfrak Z$, es decir, su cuerpo compuesto, es el cuerpo de descomposición minimal de $\mathfrak Z$.}

\subsection*{§ 9. Los isomorfismos de $Z_1$ sobre los $Z_i$}

Restringida a $\mathfrak Z$, la aplicación $\Theta_i$ es un isomorfismo de $\mathfrak Z$ sobre $Z_i$. Por consiguiente,

\[
S_i=\Theta_i\Theta_1^{-1}
\]

es un isomorfismo de $Z_1$ sobre $Z_i$. Como las $\Theta_i$ son distintas, también lo son los $S_i$. Estos son precisamente los isomorfismos de un cuerpo $Z_1$ sobre sus conjugados que se consideran habitualmente en la teoría de Galois. No existen otros, pues no existen otras aplicaciones $\Theta_i$.

\subsection*{§ 10. El grupo de Galois}

Si $\mathfrak Z$ es un cuerpo de Galois, esto es, si los $Z_i$ coinciden como conjuntos de elementos, los $S_i$ forman el grupo de todos los isomorfismos de $Z$ sobre sí mismo que fijan $\mathsf P$. En efecto, todo producto $S_iS_j$ es un automorfismo de $Z$ y, como no hay otros automorfismos aparte de los $S_i$, debe ser igual a algún $S_k$. Los $S_i$ forman, pues, el grupo de todos los automorfismos de $Z$ que fijan elemento a elemento a $\mathsf P$; este es el grupo de Galois de $Z$ respecto de $\mathsf P$.

\subsection*{§ 11. El teorema fundamental de la teoría de Galois}

Sea ahora $\mathfrak Z$ una extensión de Galois de primera especie de $\mathsf P$. Demostraremos el teorema fundamental de la teoría de Galois:

\emph{Los cuerpos $\Gamma$ intermedios entre $\mathsf P$ y $\mathfrak Z$ pueden ponerse en correspondencia biunívoca con los subgrupos $\mathfrak H$ del grupo de Galois $\mathfrak G$ de $Z$ respecto de $\mathsf P$, de tal modo que el grupo $\mathfrak H$ asociado a $\Gamma$ conste de todos los automorfismos de $Z$ que fijan elemento a elemento a $\Gamma$, y que $\Gamma$ sea el conjunto de todos los elementos de $Z$ fijos bajo todos los automorfismos de $\mathfrak H$.}

En forma abreviada:

\begin{enumerate}
\item Si $\mathfrak H$ es el grupo de invariancia de $\Gamma$, entonces $\Gamma$ es el cuerpo de invariantes de $\mathfrak H$.
\item Si $\Gamma$ es el cuerpo de invariantes de $\mathfrak H$, entonces $\mathfrak H$ es el grupo de invariancia de $\Gamma$.
\end{enumerate}

\paragraph{1. La primera afirmación se reduce al lema siguiente.}

\emph{Si $\mathsf P\subseteq\Sigma\subseteq\mathsf T\subseteq\mathsf Z$ y $\mathsf T$ tiene grado $t$ sobre $\Sigma$, todo isomorfismo de $\Sigma$ sobre un cuerpo conjugado ---necesariamente contenido en $\mathsf Z$--- admite exactamente $t$ prolongaciones distintas a isomorfismos de $\mathsf T$ sobre cuerpos conjugados.}

Supongamos, en efecto, que $\mathfrak H$ es el grupo de invariancia de $\Sigma$ y que $\mathsf T$ es el cuerpo de invariantes de $\mathfrak H$. Los elementos de $\mathfrak H$ son las prolongaciones a todo $\mathsf Z$ del isomorfismo identidad de $\Sigma$. Al restringir estas prolongaciones a $\mathsf T$, el lema muestra que se dividen en $t$ clases: los isomorfismos de una misma clase coinciden sobre $\mathsf T$, mientras que los de clases distintas son distintos sobre $\mathsf T$. Pero como $\mathsf T$ es el cuerpo de invariantes de $\mathfrak H$, ha de ser $t=1$ y $\mathsf T=\Sigma$, como se quería demostrar.

Para demostrar el lema volvemos a los cuerpos $\mathfrak S,\mathfrak T$ intermedios entre $\mathsf P$ y $\mathfrak Z$ que $\Theta_1$ aplica sobre $\Sigma$ y $\mathsf T$. Debemos probar entonces:

\emph{Todo isomorfismo $H_i$ de $\mathfrak S$ sobre un cuerpo $\Sigma_i$ puede prolongarse de exactamente $t$ maneras distintas a isomorfismos de $\mathfrak T$ sobre subcuerpos $\mathsf T_i$ de $\mathsf Z$.}

Como $\mathfrak Z$ es de primera especie sobre $\mathsf P$, el cociente $\mathfrak Z_\Omega/\mathfrak C$ coincide con $\mathfrak Z_\Omega$. Por tanto, $\mathfrak S_\Omega$ se descompone en tantos cuerpos como indica su grado $s$:

\[
\mathfrak S_\Omega=E_1\Omega+\cdots+E_s\Omega.
\]

Los $E_\nu\Omega$, considerados como módulos de representación, engendran los $s$ isomorfismos de $\mathfrak S$ sobre los $s$ subcuerpos $\Sigma_1,\ldots,\Sigma_s$. Para hallar los isomorfismos de $\mathfrak T$ necesitamos $\mathfrak T_\Omega$. Como los $E_\nu\Omega$ son ideales de $\mathfrak S_\Omega$, se tiene $E_\nu\Omega=\mathfrak S_\Omega E_\nu$. Además,

\[
\mathfrak T_\Omega=\mathfrak T\Omega=\mathfrak T\mathfrak S_\Omega.
\]

Por consiguiente,

\[
\mathfrak T_\Omega
=E_1\mathfrak S_\Omega\mathfrak T+\cdots+E_s\mathfrak S_\Omega\mathfrak T
=E_1\mathfrak T_\Omega+\cdots+E_s\mathfrak T_\Omega.
\]

Esta es una descomposición de $\mathfrak T_\Omega$ en ideales. Cada $E_\nu\mathfrak T_\Omega$ se descompone en $t$ ideales simples. En efecto, valen las relaciones de grados

\[
[\mathfrak T_\Omega E_\nu:\mathfrak S_\Omega E_\nu]
\leq[\mathfrak T_\Omega:\mathfrak S_\Omega]
=[\mathfrak T:\mathfrak S]=t;
\]

como $\mathfrak S_\Omega E_\nu=\Omega E_\nu\simeq\Omega$, resulta

\[
[\mathfrak T_\Omega E_\nu:\Omega]
=[\mathfrak T_\Omega E_\nu:\mathfrak S_\Omega E_\nu]\leq t.
\]

Pero la suma de todos los grados $[\mathfrak T_\Omega E_\nu:\Omega]$ debe ser $st$; por tanto,

\[
[\mathfrak T_\Omega E_\nu:\Omega]=t,
\qquad
E_\nu\mathfrak T_\Omega=e_\nu^{(1)}\Omega+\cdots+e_\nu^{(t)}\Omega.
\]

Cada $e_\nu^{(i)}\Omega$ proporciona una representación irreducible de $\mathfrak T_\Omega$. Solo falta demostrar que todas estas representaciones, restringidas a $\mathfrak S_\Omega$, coinciden con la engendrada por $E_\nu\Omega$. Esto es claro: de

\[
sE_\nu=E_\nu\sigma_\nu
\qquad
(s\in\mathfrak S_\Omega\text{ se representa por }\sigma_\nu)
\]

se sigue

\[
\sum_i se_\nu^{(i)}=\sum_i e_\nu^{(i)}\sigma_\nu;
\]

y, por la unicidad de la representación como suma,

\[
se_\nu^{(i)}=e_\nu^{(i)}\sigma_\nu.
\]

Así, también las nuevas representaciones asignan al elemento $s$ el elemento $\sigma_\nu$.

\paragraph{2. Sea $\mathsf T$ el cuerpo de invariantes de $\mathfrak H$.}

Queremos reconocer $\mathfrak H$ como grupo de invariancia de $\mathsf T$. Para ello basta exhibir un elemento de $\mathsf T$ que no permanezca fijo bajo un automorfismo que no pertenezca a $\mathfrak H$.

Los elementos $S_i$ del grupo de Galois $\mathfrak G$ de $\mathsf Z$ son los automorfismos de $\mathsf Z$ que fijan $\mathsf P$. Los hacemos actuar como automorfismos de $\mathfrak Z$ mediante la convención

\[
S_i z=z\qquad\text{si }z\text{ pertenece a }\mathfrak Z.
\]

La aplicación de $S_i$ a

\[
\mathfrak Z_{\mathsf Z}=e_1\mathsf Z+\cdots+e_n\mathsf Z
\]

debe permutar los sumandos directos $e_\nu\mathsf Z$, pues estos están determinados de manera única. En particular,

\[
S_ie_\nu=e_{\nu'},\qquad S_ie_\nu\ne S_ke_\nu\quad(i\ne k).
\]

Sean $S_1,\ldots,S_h$ los elementos de $\mathfrak H$. Elegimos una componente cualquiera $e'$, y consideramos todas las componentes obtenidas de ella por los $S_i$, $i\leq h$:

\[
e_1=S_1e',\ldots,e_h=S_he'.
\]

Si $i,j\leq h$, entonces

\[
S_ie_j=S_iS_je'=S_ke'=e_k\qquad\text{para algún }k\leq h.
\]

Por tanto, todo $S_i\in\mathfrak H$ permuta entre sí $e_1,\ldots,e_h$, y

\[
E_1=e_1+\cdots+e_h
\]

es un elemento de $\mathfrak Z$ invariante bajo todos los $S_i\in\mathfrak H$. Al expresar $E_1$ respecto de una base de $\mathfrak Z$, sus coeficientes deben ser elementos de $\mathsf Z$ invariantes bajo $\mathfrak H$; pertenecen, pues, a $\mathsf T$.

El elemento $E_1$ no es invariante bajo ningún $S_i\notin\mathfrak H$. En efecto, para tal $S_i$, la suma

\[
S_iE_1=e_i+\cdots
\]

contiene la componente $e_i$, que no figura entre $e_1,\ldots,e_h$; pero la representación de $E_1$ como suma es única.

\subsection*{§ 12. Significado formal de las componentes $e_i$}

Fijemos una base $z_1,\ldots,z_n$ de $\mathfrak Z$:

\[
\mathfrak Z=z_1\mathsf P+\cdots+z_n\mathsf P.
\]

Ahora solo suponemos que $\mathfrak Z$ es de primera especie. El isomorfismo $\Theta_i$ aplica el sistema $z_1,\ldots,z_n$ sobre una base $\zeta_1^{(i)},\ldots,\zeta_n^{(i)}$ de $\mathsf Z_i$ mediante las relaciones

\[
z_\nu e_i=e_i\zeta_\nu^{(i)}.
\]

Recíprocamente, la componente $e_j$ de la unidad debe poder expresarse mediante los $z_\nu$ con coeficientes en $\mathsf Z$:

\[
e_j=\sum_\nu\varrho_\nu^{(j)}z_\nu.
\]

Por tanto,

\[
e_j=e_j^2=\sum_\nu\varrho_\nu^{(j)}z_\nu e_j
=\sum_\nu\varrho_\nu^{(j)}\zeta_\nu^{(j)}e_j,
\]

y

\[
0=e_je_k=\sum_\nu\varrho_\nu^{(j)}z_\nu e_k
=\sum_\nu\varrho_\nu^{(j)}\zeta_\nu^{(k)}e_k.
\]

Así,

\[
\sum_\nu\varrho_\nu^{(j)}\zeta_\nu^{(j)}=1,
\qquad
\sum_\nu\varrho_\nu^{(j)}\zeta_\nu^{(k)}=0\quad(j\ne k).
\]

Esto significa que las matrices

\[
\begin{pmatrix}
\zeta_1^{(1)}&\cdots&\zeta_n^{(1)}\\
\vdots&&\vdots\\
\zeta_1^{(n)}&\cdots&\zeta_n^{(n)}
\end{pmatrix},
\qquad
\begin{pmatrix}
\varrho_1^{(1)}&\cdots&\varrho_1^{(n)}\\
\vdots&&\vdots\\
\varrho_n^{(1)}&\cdots&\varrho_n^{(n)}
\end{pmatrix}
\]

son inversas entre sí. Esto también se expresa diciendo que $\varrho_1^{(i)},\ldots,\varrho_n^{(i)}$ es la base complementaria de $\zeta_1^{(i)},\ldots,\zeta_n^{(i)}$. Esta base interviene en la teoría de la diferente de un cuerpo numérico.

\subsection*{§ 13. El teorema general de prolongación}

El teorema demostrado anteriormente sobre la prolongación de los $H_i$ admite la generalización siguiente.

Sea $\mathfrak Z=z_1\mathsf P+\cdots+z_n\mathsf P$ un sistema conmutativo, hipercomplejo y completamente reducible, y sean $\mathfrak S$ y $\mathfrak T$ subsistemas ---necesariamente también completamente reducibles--- tales que

\[
\mathsf P\subseteq\mathfrak S\subseteq\mathfrak T\subseteq\mathfrak Z.
\]

Según M.Z., §13, cada uno de los tres sistemas posee unidad, aunque estas unidades no tienen por qué coincidir. En particular, no coinciden cuando $\mathfrak S$ y $\mathfrak T$ son ideales de $\mathfrak Z$ distintos de $\mathfrak Z$. Sea $E$ la unidad de $\mathfrak S$. El sistema $\mathfrak T E$ es un módulo de rango finito sobre $\mathfrak S$:

\[
E\mathfrak T=a_1\mathfrak S+\cdots+a_t\mathfrak S,
\]

aunque $\mathfrak T$ en general no lo sea.

\emph{Demostración.} El sistema $\mathfrak S$ es suma directa de cuerpos,

\[
\mathfrak S=\mathfrak R_1+\cdots+\mathfrak R_r.
\]

A ella corresponde una descomposición en suma directa de $\mathfrak T E$. Cada sumando de $\mathfrak T E$ debe contener alguno de los cuerpos $\mathfrak R_i$; de lo contrario sería anulado por $E$, siendo $E$ la unidad de $\mathfrak T E$. Por tanto $\mathfrak T E$ es suma directa de anillos de extensión de los $\mathfrak R_i$, esto es,

\[
\mathfrak T E=a_{11}\mathfrak R_1+\cdots+a_{1p_1}\mathfrak R_1
+\cdots+a_{rp_r}\mathfrak R_r,
\]

o bien $E\mathfrak T=a_{11}\mathfrak S+\cdots+a_{rp_r}\mathfrak S$, como se quería demostrar.

\begin{center}
Teorema general de prolongación
\end{center}

\emph{Toda representación irreducible de $\mathfrak S$ en $\Omega$, dada por un homomorfismo $H_i$, puede prolongarse de exactamente $t$ maneras distintas a representaciones de $\mathfrak T$.}

La demostración es análoga a la anterior. Se tiene

\[
\mathfrak S_\Omega=E_1\Omega+\cdots+E_s\Omega,
\qquad E_\nu\Omega=\mathfrak S_\Omega E_\nu,
\]

\[
\mathfrak T_\Omega E=\mathfrak T E\Omega
=\mathfrak T\mathfrak S\Omega=\mathfrak T\mathfrak S_\Omega,
\]

\[
\mathfrak T_\Omega E
=\mathfrak T(E_1\mathfrak S_\Omega+\cdots+E_s\mathfrak S_\Omega)
=E_1\mathfrak T_\Omega+\cdots+E_s\mathfrak T_\Omega,
\]

y

\[
\begin{aligned}
[\mathfrak T_\Omega E_\nu:\Omega]
&=[\mathfrak T_\Omega E_\nu:\Omega E_\nu]
=[\mathfrak T_\Omega E_\nu:E_\nu\mathfrak S_\Omega]\\
&\leq[\mathfrak T_\Omega E:\mathfrak S_\Omega]
=[\mathfrak T E:\mathfrak S]=t.
\end{aligned}
\]

Como

\[
\sum_\nu[\mathfrak T_\Omega E_\nu:\Omega]=st,
\]

resulta

\[
[\mathfrak T_\Omega E_\nu:\Omega]=t.
\]

Por consiguiente, $\mathfrak T_\Omega E_\nu$ se descompone en $t$ ideales. Cada uno de ellos, como sumando de $\mathfrak T_\Omega E$ y por tanto también de $\mathfrak T_\Omega$, proporciona una representación de $\mathfrak T_\Omega$ que, por el mismo argumento anterior, es una prolongación de la representación dada por $H_\nu$. Las demás representaciones irreducibles de $\mathfrak T_\Omega$ ---engendradas por los sumandos de $\mathfrak T_\Omega$ que no están contenidos en $\mathfrak T_\Omega E_\nu$--- no pueden prolongar las representaciones de $\mathfrak S_\Omega$ engendradas por los $H_\nu$, pues esos ideales son anulados por $\mathfrak T_\Omega E_\nu$ y, con mayor razón, por $\mathfrak S_\Omega$.

\section*{Capítulo III. Grupos abelianos}
